%% Title and Preprocessing
\documentclass[titlepage]{article}\pagestyle{empty}
\usepackage{amsmath}

\author{Kevin Fang}
\title{Homework \#11}
\date{\today}
\usepackage[left=1cm, right=2cm, vmargin=1cm]{geometry}
\usepackage{amsfonts}

\def\therefore{\boldsymbol{\text{ }
\leavevmode
\lower0.4ex\hbox{$\cdot$}
\kern-.5em\raise0.7ex\hbox{$\cdot$}
\kern-0.55em\lower0.4ex\hbox{$\cdot$}
\thinspace\text{ }}}

\usepackage{mathtools}
\DeclarePairedDelimiter\ceil{\lceil}{\rceil}
\DeclarePairedDelimiter\floor{\lfloor}{\rfloor}

% Start Document Here
\begin{document}
\maketitle

\pagebreak
\section*{Question 5.}
\subsection*{a.}
Base case: $n=1$.\\
For $n=1$,\\
$n^3+2n=1^3+2*1=3$\\
Thus, $3$ divides $n^3+2n$ for $n=1$.\\
Induction: For all $k \geq 1$, if $3$ divides $k^2 + 2k$, we prove 3 divides $(k+1)^2+2(k+1)$.\\
$k^3+2k=3m$,$m \in \mathbb{Z}^+$\\
$(k+1)^3+2(k+1)=k^3+3k^2+3k+1+2k+2$\\
$=(k^3+2k)+(3k^2+3k+3)$\\
$=3m+3(k^2+k+1)$\\
$=3(m+k^2+2k+1)$\\
Since $m+k^2+2k+1$ is an positive integer, 3 divides $(k+1)^2+2(k+1)$.\\

\subsection*{b.}
Base case: $n=2$\\
For $n=2$, $2$ is a prime.\\
Induction: Assume for all $k \geq 2$, all integers can be written as a product of prime.\\
We then prove $k+1$ can be written as a product of prime.\\
If $k+1$ is prime, then it can be written as $1*(k+1)$\\
If $k+1$ is not prime, then it can be written as\\
$k+1=a*b$ for $2 \leq a, b \leq k$\\
Since we assumed for all $k \geq 2$, all integers can be written as a product of prime.\\
Thus, $a$ and $b$ are products of prime.\\
Thus, $k+1$ is a product of prime.\\

\pagebreak
\section*{Question 6.}
\subsection*{(a): 7.4.1}
\subsubsection*{A.} 
$P(3)$ is true since $1^2 + 2^2+3^2 = (3(3+1)((3 * 2) + 1))/6 = 14$.
\subsubsection*{B.}
$P(k)=\sum_{j=1}^{k}j^2=\frac{k(k+1)(2k+1)}{6}$
\subsubsection*{C.}
$P(k+1)=\sum_{j=1}^{k+1}j^2=\frac{(k+1)(k+1+1)(2(k+1)+1)}{6}$
\subsubsection*{D.}
$P(1) = \frac{6}{6} = 1$. Thus $P(1)$ is true.
\subsubsection*{E.}
We prove in the inductive step that: $\sum_{j=1}^{k+1}j^2=\frac{(k+1)(k+1+1)(2(k+1)+1)}{6}$\\
For all $k \geq 1$, $P(k)$ implies $P(k+1)$.
\subsubsection*{F.}
The inductive hypothesis is $P(k)=\sum_{j=1}^{k}j^2=\frac{(k+1)(2k+1)}{6}$
\subsubsection*{G.}
Base case: $n=1$\\
$\sum_{j=1}^{1}k^2=1^2=1$\\
$1*(1+1)*(2*1+1)/6=(1*2*3)/6=1$\\
Thus, $P(1)$ is true.\\
Induction: For all $k \geq 1$, if $P(k)$ is true, we prove $P(k+1)$ is also true.\\
$\sum_{j=1}^{k+1}(k+1)^2=\sum_{j=1}^{k}k^2+(k+1)^2$\\
$=\frac{k(k+1)(2k+1)}{6}+(k+1)^2$\\
$=\frac{k(k+1)(2k+1)}{6}+\frac{6(k+1)(k+1)}{6}$\\
$=\frac{(k+1)(2k^2+7k+6)}{6}$\\
$=\frac{(k+1)(k+2)(2k+3)}{6}$\\
$=\frac{(k+1)((k+1)+1)(2(k+1)+1)}{6}$\\
$=\frac{l(l+1)(2l+1)}{6}$ for $l=k+1$\\
Thus, for all $k\geq1$, $P(k+1)$ is true if $P(k)$ is true.\\
\pagebreak
\subsection*{(b):7.4.3}
Prove that for $n \geq 1$: $\sum_{j=1}^{n}\frac{1}{j^2}\le2-\frac{1}{n}$.\\
Base case: $n = 1$:\\
\[f(1)=\frac{1}{1^2}\le2-\frac{1}{1}=1\leq 1\]\\
Base case is true.\\
Induction:\\
We assume \[\sum_{j=1}^{k}\frac{1}{j^2}\le2-\frac{1}{k}\]
We prove that:
\[\sum_{j=1}^{k+1}\frac{1}{j^2}\le2-\frac{1}{k+1}\]
$\sum_{j=1}^{k+1}\frac{1}{(k+1)^2}=\sum_{j=1}^{k}\frac{1}{k^2}+\frac{1}{(k+1)^2}$\\
$\leq2-\frac{1}{k}+\frac{1}{(k+1)^2}$\\
$\leq2-\frac{1}{k}+\frac{1}{k(k+1)}$\\
$\leq2-\frac{(k+1)}{k(k+1)}+\frac{1}{k(k+1)}$\\
$\leq2-\frac{k}{k(k+1)}$\\
$\leq2-\frac{1}{(k+1)}$\\
Thus, for all $k\geq1$, $P(k+1)$ is true if $P(k)$ is true.\\
\subsection*{(c):7.5.1}
Base case: $n=1$\\
For $n=1$, $3^{2*1}-1=8=2*4$\\
Thus, $P(1)$ is true.\\
Induction: For all $k\geq1$, we prove $P(k+1)$ is true if $P(k)$ is true.\\
$3^{2k}-1=4m, m \in \mathbb{Z}^+$\\
$3^{2(k+1)}-1=3^{2k+2}-1$\\
$=3^2*3^{2k}-1$\\
$=9*3^{2k}-9+8$\\
$=9*(3^{2k}-1)+8$\\
$=9*4m+8$\\
$=4*(9m+2)$\\
Since $9m+2$ is a positive integer, $4$ evenly divides $3^{2(k+1)}-1$. Thus, for all $k\geq1$, $P(k+1)$ is true if $P(k)$ is true.
\end{document}